
\section{怎样辨析同义词}

同学们在作文里，有时会写出这类句子：“勇敢的人民解放军是一支战无不胜的队伍。”老师往往把“勇敢”改成了“英勇”。老师为什么要这样改呢？你看，这两个词都有“有胆量、不怕危险和困难”的意思，而“英勇”一词更有“勇敢出众或特别勇敢”的意思。很明显，在前面所说的句子里，只有用“英勇”这个词才能恰当地表达出人民解放军的英雄本色。可见，只有掌握了同义词的意义和用法，才能从一组同义词里选取最恰当的词来准确地反映客观事物和表达思想感情。

那么，怎样才能准确地辨析和使用同义词呢？

我们知道，同义词是指语言里意义相同或者相近的词。但是，同义词“同中有异”，这“异”主要表现在词的意义、感情色彩和具体用法等几个方面的细微差别上。因此，要想准确辨析同义词，一定要努力找出这“异”来。具体地说，至少应注意以下几点：

\textbf{1. 注意区别词义的使用范围。}

我们看，“凤凰岭战斗，是解放战争时期孟良崮战役中的一次激烈的战斗。”在这句话中，“战争”的范围最大，包括民族与民族、国家与国家、阶级与阶级或政治集团与政治集团之间的武装斗争；而“战役”只是在一定时间内一系列战斗的总和，范围较小，至于“战斗”呢，是指敌对双方所进行的武装冲突，范围最小。它们之间当然不能互相替换。

一般说来，同义词使用范围的大或小，可以从同义词的不同词素上看出来。比如“支援”和“声援”都有支持、援助的意思，但包括的范围不同。“支援”包括舆论、物质和行动方面的支持和帮助，“声援”则偏重于舆论和道义上的支持。“支援”所包括的范围较“声援”为广。再如“功劳”和“功勋”都有对事业作出贡献的意思，但“功劳”包括大大小小的贡献，范围大，而“功勋”特指对国家、人民建立特别重大的功劳，范围就比“功劳”小。（当然，“功勋”又比“功劳”的程度重）。其他如“边界、边境”、“年月、年代”，“品质、性质”，“房子、房屋”，“事故、事件”等等，都可以用这种方法辨析出前者的范围比后者小。

我们根据这种辨析方法，就可以避免或者改正由于不明确同义词词义范围的大小而造成的病句。例如：

（1）这次黑板报的编写工作，你出了不少力，功勋（应为“功劳”）不小。

（2）那时房东的四间房屋（应为“房子”）就被匪兵强占了三间。

（3）我们班的文娱（应为“文体”）活动开展得很好，经常大唱革命歌曲，并定期举行篮球比赛。

（4）早上气候（应为“天气”）很好，没有多久却阴云密布，下起雨来。

\textbf{2. 注意区别词义的轻重程度。}

这一类在词义轻重程度上有细微差别的同义词，以动词和形容词居多数。

我们看，“他成绩优良”，“他成绩优秀”，“他成绩优异”，这三个句子中的“优良”，“优秀”，“优异”三个同义词，都是形容词，都有好的意思。但是，“优良”是指超出一般的好，“优秀”是指非常好，“优异”则是出类拔萃的好。前面的词义换次比后面的轻。

我们再看，“你应劝止他”，“你应阻止他”，“你应制止他”，这三个句子中的三个同义词都是动词，都有“使不能进行”的意思。但是，“劝止”指用道理劝说人停止行动，“阻止”指阻拦人去干某事，而“制止”则是强迫人停止，不允许继续活动。前面的词义换次比后面的轻。

这一类同义词大多数可以从它们不同词素的含义上辨析出细微的差别来。比如“请求、恳求、哀求”，“喜爱、热爱、酷爱”、“轻视、蔑视、鄙视”等等，都可以用这种方法辨析出前者的程度换次比后面的轻。

我们根据这种辨析方法，就可以避免或者改正由于不明确同义词词义的轻重程度而造成的病句。例如：

（1）有些人开会常常早退，这种恶劣（应为“不良”）现象应该受到批评。

（2）革命干部要教育好自己的子女，至少不能纵容（应为“容忍”）子女去干坏事。

（3）今天，我仔细攻读（应为“阅读”）了《木偶奇遇记》，很受教育。

（4）帝国主义的侵略政策妨碍（应为“妨害”）了世界和平。

\textbf{3. 注意区别词义的褒贬倾向。}

有些词表示的意思虽然相同或者相近，但它们的感情色彩并不一样。比如“鼓励”和“怂恿”，都有鼓动人做某件事的意思。但是，“鼓励”是勉励人做好事，是褒义词；“怂恿”是煽动人作坏事，是贬义词。至于“鼓动”，有时可以用在好的方面，如“我们一定要做好宣传鼓动工作”，这时它与“鼓励”的意思有点接近；有时可以用在坏的方向，如“反革命分子暗中鼓动闹事”。这时它与“怂恿”的意思较为接近，说明它是个中性词。

辨析这类同义词比较好办，只要把它们分别放到表示褒贬的句子中去，立即可以检验出来。比如“团结、勾结、联合”，“成果、后果、结果”，“果断、武断、判断”，“宏大、庞大、巨大”等等，它们中的第一个词只能用在好的方面，第二个词只能用在坏的方面，第三个词有时可以用在好的方面，有时可以用在坏的方面。我们在写文章选词语时，如果对某个词的褒贬色彩把握不定的话，不妨再从反面造个句子，就可以确定这个词的褒贬倾向用在这儿是否妥当了。

我们根据这种辨析方法，就可以避免或者改正由于不明确同义词词义的感情色彩而造成的病句。例如：

（1）大家对小王的批评正确而尖刻（应为“尖锐”）。

（2）这些信都连篇累牍（应为“自始至终”）地充满了他对老师的深情。

（3）“四人帮”大力（应为“大肆”）污蔑和攻击社会主义制度，妄图复辟资本主义。

（4）这篇文章的主要企图（应为“意图”），并不在于赞美自然景色，而是在于歌颂祖国的社会主义建设。

\textbf{4. 注意区别词义的语体色彩。}

词从语体上讲，有书面语与口语、方言与普通话、日常用语与专门用语之别，有些词常在某种范围内使用，形成各自的语体色彩。这一类的同义词，我们应从它们的语体色彩的不同去分辨。

例如：“祖国啊，我的母亲！”、“今天是周恩来总理的八十九周年诞辰纪念日。”“彻底的唯物主义者是无所畏惧的。” 这儿的“母亲”、“诞辰”、“畏惧”，是书面语，较庄重；而“妈妈”、“生日”、“害怕”，是口语，较通俗，没有庄重的色彩，它们之间是不能替换的。其它如：“恫吓”与“吓唬”，“吝啬”与“小气”，“黎明”与“天亮”，“措施”与“办法”，“部署”与“安排”，“聆听”与“听”，“铭记”与“记”，“寒冷”与“冷”，“观看”与“看”，等等，前者多用于书面语，后者多用于口语。

再如：“你说啥？”“他们才是真正的里手。”这儿的\\
“啥”，“里手”，都是方言词，带有地方色彩；而“什么”，\\
“内行”，是普通话，取消了地方色彩。在特定的表达地方色彩的语言环境里，二者之间也不能随意替换。

又如：“率领”和“带领”，“会见”和“见”，“给予”和“给”，“陪同”和“陪”，它们中的前者都含有庄重的色彩，大多用于公文、贺电、新闻报道等文体，后者则用于一般的文章。“氧化钙”和“石灰”，“玉蜀黍”和“玉米”等，词义全同，只是它们中的前者是科学术语，后者是日常用语，在特定的文体里一般不能混用。至于“飞翔、翱翔”和“飞”，“心灵、心田、心扉、心弦”和“心”，“寂静、宁静、静穆、静谧”和“静”，等等，前者这些双音词大多用于文艺作品，后者这些单音词则常用于一般的文章。
这一类同义词，我们只要弄清它们的语体色彩和语言环境，就不会将它们混淆起来而误用的。

我们根据这种辨析方法，就可以避免或者改正由于不明确同义词词义的语体色彩而造成的病句。例如：

（1）妹妹才十岁，妈妈就给她过诞辰（应为“生日”），我都十几岁了，为什么不给过我呢？

（2）我们来到红光造纸厂参观，张厂长亲切地接见（应为“接待”）了我们。

（3）爷爷的下肢（应为“腿”）不灵便，我就搀着他去看电影。

（4）为人处事，首先心灵（应为“心”）要真。

\textbf{5. 注意区别与其他词语的搭配关系。}

一组同义词中，一个词跟哪些词或词组相搭配，是有一定限制的。有些词经常同某些词配合使用，形成了比较固定的配合习惯。

例如：“改进”和“改善”，这两个词都是动词，都有“改变原有情况向好的方向发展”的意思，但是“改进”着重在“使有所进步”，“改善”着重在“使变得好些”，所以跟它们相搭配的词也不相同。“改进”一般与“工作、方法、技术”等词相搭配，“改善”一般与“生活、条件、关系”等词相搭配。如果把二者相搭配的词用颠倒了，就不符合已经形成的习惯，而成为一种搭配不当的语病了。

再如：

$\begin{cases}
    \text{表现 --- 品质、生活、思想、感情、作风}\\
    \text{体现 --- 原则、要求、理想、精神、关系}
\end{cases}$

$\begin{cases}
    \text{发扬 --- 优点、成绩、传统、作风、精神}\\
    \text{发挥 --- 智慧、干劲、作用、威力、创造性}
\end{cases}$

$\begin{cases}
    \text{履行 --- 诺言、合同、手续、条约、义务}\\
    \text{执行 --- 命令、任务、方针、路线、政策}
\end{cases}$

$\begin{cases}
    \text{转移 --- 视线、阵地、目标、方向、火力点}\\
    \text{转变 --- 立场、态度、作风、思想、立足点}
\end{cases}$

$\begin{cases}
    \text{交换 --- 意见、资料、礼物、人质、情报}\\
    \text{交流 --- 思想、经验、物资、人才、信息}
\end{cases}$

我们根据这种辨析方法，就可以避免或者改正由于不明确同义词与其他词语相搭配的关系而造成的病句。例如：

（1）我们要发挥（应为“发扬”）艰苦奋斗的精神，努力为四化建设多作贡献。

（2）每个人都要知法守法，执行（应为“履行”）自己应尽的社会义务。

（3）全班同学交换（应为“交流”）了学习经验之后，都感到学习效率明显提高了。

（4）经过班主任的耐心教育，我不再对自己所担负（应为“担任”）的班长职务推辞了。

总的说来，我们在辨析同义词时，应该把同义词放在句子中，联系上下文，综合运用以上所说的各点，找出同义词在词义、色彩和用法方面的不同，着重揭示它们之间的主要差别。如果能有意识地作一些比较和辨析同义词的练习，把掌握的有关知识化为纯熟运用的能力，我们自能区别和使用同义词。例如：

“他尽管浑身被敌人打得（ ），仍然英勇不屈。”（遍体鳞伤，体无完肤，皮开肉绽）

句末括号里的三个词语，都有“伤势非常严重”的意思，如果要求我们选择一个恰当的词语填在句中的括号内，

那么，究竟填哪一个为好呢？这只要联系前面的“浑身”二字加以思考，就能作出正确的选择。
\newpage